\documentclass[12pt]{article}
\input{iati_structure.tex}
\renewcommand{\bottomtitlespace}{2cm}

\usepackage[english]{babel}

\begin{document}

\renewcommand{\headerLang}{English}
\renewcommand{\problemName}{AB23. TRIANGLE}

\renewcommand{\headerLeft}{IATI Day 2, Senior group}
\renewcommand{\headerRightFirst}{XVII INTERNATIONAL ADVANCED TOURNAMENT IN INFORMATICS}
\renewcommand{\headerRightSecond}{ONLINE 2026}

\renewcommand{\tl}{$10$ sec.}
\renewcommand{\ml}{$1024$ MB}
\problem{Problem \problemName}
\addauthor{Author: Boris Mihov}{2026}{04}{19}

Novo rukovodstvo državne komisije, koju predstavljaju koordinatori grupa A i B, želi da održi striktan plan i program datih zadataka za biranje reprezentacije. U tu svrhu, moraće da daju zadatak iz geometrije. Štaviše -- broj interaktivnih zadataka datih do sad je ispod standarda. Da bi rešio krizu u izboru tima za \texttt{IOI}, Vasa je odlučio da da zadatak koji je istovremeno interaktivan i geometrijski. Ima sledeću postavku:

Vasa ima skrivenu permutaciju $P_0,P_1,...,P_{N-1}$ niza $\{1,2,3,...,N\}$. Tvoje je da pogodiš tu permutaciju. U tu svrhu, možeš postaviti Vasi sledeće pitanje: ``Da li je moguće formirati trougao pozitivne površine, dužina stranica $P_A,P_B,P_C$?''

Primerimo da je poznato (po nejednakosti trougla) da je to moguće ako i samo ako je:

\[P_A + P_B > P_C,\ P_A + P_C > P_B \quad \text{and} \quad P_B + P_C > P_A.\]

Napiši program \textbf{\texttt{triangle}}, koji sadrži funkciju \texttt{solve}, koja će se compile-ovati zajedno sa Vasinim programom i koja će s njim komunicirati tako što će mu postavljati pitanja oblika koji je naveden iznad. Na kraju izvršavanja, mora pogoditi permutaciju.

\subsection{Detalji implementacije}
Potrebno je implementirati \texttt{solve} funkciju:
\begin{lstlisting}
 std::vector<int> solve(solve N)
\end{lstlisting}
\begin{itemize}
	\item $N$: dužina permutacije.
\end{itemize}

Ova funkcija će biti pozvana $T$ puta po test primeru -- jednom po podtestu, pri čemu će svaki imati istu vrednost $N$ i treba da vrati skrivenu permutaciju za podtest. Tvoj program može pozvati Vasinu funkciju \texttt{query}:
\begin{lstlisting}
 bool query(int A, int B, int C)
\end{lstlisting}
\begin{itemize}
	\item $A$, $B$, $C$: indeksi stranica trougla $P_A, P_B, P_C$.
\end{itemize}

Funkcija vraća \texttt{true}, ako trougao pozitivne površine dužna stranica $P_A,P_B,P_C$ postoji, a \texttt{false} u suprotnom.

\subsection{Ograničenja}
\vspace{0.1em}
\begin{itemize}
	\item $N = T = 1~000$
\end{itemize}

\subsection{Podzadaci}
\begin{table}[H]
	\begin{tblr}{width=\linewidth,colspec={|Q[c,m]|Q[c,m]|X[c,m]|}}
		\hline
		\textbf{Podzadatak} & \textbf{Poeni} & \textbf{Opis}\\
		\hline
		$1$ & $100$ & Podzadatak se sastoji od $1$ test primera sa $T=1~000$ nasumično generisanih permutacija, gde svaka sadrži $N=1~000$ elemenata.  \\ 
		\hline
	\end{tblr}
	\caption*{The points for a given subtask are obtained only if all the tests for it.}
\end{table}
\FloatBarrier

\newpage
\subsection{Bodovanje}

Neka je $Q_{contestant}$ prosečan broj upita koji tvoj program koristi za 1 poziv funkcije \texttt{solve} na jednom podtestu, i neka je $Q_{target}=8770$. Broj poena na zadatku je jednak:

\[
\begin{cases}
0 & \text{if } Q_{contestant} > 2\times 10^6, \\
100 & \text{if } Q_{contestant} \le Q_{target}, \\
100 \times \max\left(0.15,\; 1-\sqrt{1-\left(\frac{Q_{target}}{Q_{contestant}}\right)^{0.55}}\right) & \text{if } Q_{contestant} > Q_{target}.
\end{cases}
\]

\subsection{Sample interaction}
For this sample interaction $N=4$, $T=2$.
\begin{table}[H]
	\begin{tblr}{|c|c|}
		\hline
		\textbf{Contestant action} & \textbf{Jury action} \\
		\hline
		& \texttt{solve(4)} \\
		\hline
		\texttt{triangle(0, 0, 0)} & \texttt{true} \\
		\hline
        \texttt{triangle(0, 1, 2)} & \texttt{false} \\
		\hline
        \texttt{triangle(0, 1, 3)} & \texttt{false} \\
		\hline
        \texttt{triangle(0, 2, 3)} & \texttt{true} \\
		\hline
        \texttt{triangle(1, 2, 3)} & \texttt{false} \\
		\hline
        \texttt{return \{3, 1, 2, 4\};} & \\
        \hline
        & \texttt{solve(4)} \\
		\hline
		\texttt{triangle(0, 1, 2)} & \texttt{false} \\
		\hline
        \texttt{triangle(0, 1, 3)} & \texttt{true} \\
		\hline
        \texttt{triangle(0, 2, 3)} & \texttt{false} \\
		\hline
        \texttt{triangle(1, 2, 3)} & \texttt{talse} \\
		\hline
        \texttt{return \{4, 2, 1, 3\};} & \\
		\hline
	\end{tblr}
\end{table}
\FloatBarrier
	
\subsection{Lokalno testiranje}
Input format:
\begin{itemize}
	\item linija $1$: tri cela broja $T$, $N$, i $R$ -- dužina permutacija, broj podtestova i režim izvršavanja. Za $R = 1$, lokalni grader će generisati nasumične permutacije i očekivaće broj u prvoj liniji -- $S$, što je seed za random generator. Za $R = 2$, ulaz se nastavlja:
	\item linije od $2$ do $(1+T)$: permutacije niza $\{1,2,\dots,N\}$.
\end{itemize}

Output format:
\begin{itemize}
	\item linija $1$: poruka o grešci ili prosečan broj upita za podtestove ako su sve permutacije pogođene.
\end{itemize}

\end{document}